\section{AI Agents That Matter}
\paperhead{AI Agents That Matter}{Sayash Kapoor, Benedikt Stroebl, Zachary S. Siegel, Nitya Nadgir, Arvind Narayanan}{Transactions on Machine Learning Research (TMLR), 2025}{https://arxiv.org/abs/2407.01502}{arXiv API id\_list=2407.01502, title match 1.00}

\subsection{What the paper does}

Agent research is steered by leaderboards that report accuracy alone, which the authors argue is a degenerate objective for agents specifically: the underlying model is stochastic and many agent environments return a cheap correctness signal such as a unit test, so accuracy can be purchased by sampling again. An agent that retries enough times tops the board while contributing nothing scientific, and the field then credits the gain to whatever architecture was wrapped around the retries.

From that they build a four-part indictment: accuracy is reported without cost; model evaluation (a scientific question about architectures and data) is conflated with downstream evaluation (a procurement question about dollars); benchmarks lack held-out sets at the right level of generality, so agents can hardcode their way to the top; and harnesses are non-standard enough that published numbers do not reproduce. The remedy is a two-dimensional accuracy versus dollar-cost Pareto plane, with cost as an optimization target rather than a footnote.

\subsection{How the method works}

Three instruments. (i) \textit{Cost-controlled re-evaluation.} They re-run three published HumanEval agents that share code (LDB, LATS, Reflexion) five times each, reporting mean accuracy against mean total dollar cost, using the LDB variant of HumanEval because the original supplies example tests for only 161 of its 164 problems. Against these they place three deliberately trivial baselines: \textit{Retry}, re-invoking the model at temperature zero up to five times whenever the example tests fail (viable because models are not deterministic even at temperature zero); \textit{Warming}, the same loop with temperature ramped from 0 to 0.5; and \textit{Escalation}, starting on Llama-3 8B and climbing to GPT-3.5, Llama-3 70B and GPT-4 only on a test failure. The frontier is drawn convex, since two agents can always be interpolated by invoking one with probability $p$ and the other with $1-p$.

(ii) \textit{Joint optimization.} They separate fixed cost (one-time tuning of prompt and hyperparameters) from variable cost (tokens per run), then modify DSPy on HotPotQA multi-hop retrieval and drive an Optuna search over per-module temperature, the number of few-shot examples, which examples are selected, and whether formatting instructions are included, maximizing accuracy while minimizing prompt tokens. Tuning uses 100 training samples split evenly into candidate collection and validation; 200 evaluation samples score the result.

(iii) \textit{A holdout ladder.} Four levels of agent generality, distribution-specific, task-specific, domain-general and fully general, each with a prescribed holdout (in-distribution samples, out-of-distribution samples, held-out tasks, held-out domains), applied as an audit over 17 agent benchmarks.

\subsection{What it reports}

On HumanEval the trivial baselines Pareto-dominate the published architectures. There is no significant accuracy difference between warming and the best-performing agent architecture, while for substantially similar accuracy cost differs by almost two orders of magnitude: Reflexion and LDB cost over 50\% more than warming, LATS over 50 times more, and escalation strictly improves accuracy over LDB (GPT-3.5) at less than half its cost. The authors conclude there is no evidence that ``System 2'' machinery (planning, reflection, debugging) produces the gains on this benchmark. Joint optimization on HotPotQA cuts variable cost by 53\% for GPT-3.5 and 41\% for Llama-3-70B at similar accuracy, the fixed tuning cost amortizing after roughly 1,350 tasks. NovelQA shows a protocol distorting a design decision: batching every question about a novel misrepresents sequential real use, where RAG is roughly as accurate as long context at more than 20 times less cost, whereas NovelQA shows it costing only half as much. The audit finds appropriate holdouts in 1 of 8 domain-general and 0 of 2 fully general benchmarks; WebArena's top agent STeP reaches 35.8\% accuracy partly by hardcoding policies such as appending \texttt{/user/user\_name} to the base URL, which breaks under site drift. On reproducibility they report order dependence from Reddit rate limits, published scores exceeding the maximum of their own five runs, and both LATS and STeP marking some incorrect tasks correct while dropping 1 and 8 tasks.

\subsection{How Mytheus uses this}
\begin{takeaway}
Report the five-lever cost cascade as a Pareto curve, not as the headline ``1,206 versus 8,640 calls, 86\% reduction''. That figure is exactly the uninterpretable claim this paper attacks: it says nothing about the accuracy paid for it. Build the sweep by holding the golden set fixed while varying the Haiku triage pre-gate threshold (currently about 35\% pass) and the LiteLLM tier routing policy, plotting DecisionAgent and ActionAgent extraction F1 and PersonalListenerAgent relevance precision/recall against dollars per meeting, and publish token counts alongside dollars so the curve survives price changes. Add the paper's \textit{Escalation} baseline as a real ablation: Haiku first, escalating to Sonnet only when confidence falls in the 0.4 to 0.7 band, which turns the existing gates into the escalation trigger for free. The holdout critique lands hardest on the golden-set CI harness: if prompt templates are iterated until the golden set passes and prompt changes are then gated on that same suite, the golden set has become training data. Split it into a tuning half and a sealed half never touched during iteration, and treat the 48-hour canary on 1\% of tenants as the real out-of-distribution test, holding out whole tenants and whole meeting types rather than random windows from meetings already tuned on. What does not transfer: the winning baselines all buy accuracy with extra serial model calls. Retry and warming cannot sit inside ContextGapAgent's p95 1.5 s peer round trip or the p95 4 s agent output budget, so Mytheus needs a third axis, p95 latency, and must confine retry-style gains to post-meeting synthesis. Conversely the paper's admitted human-in-the-loop gap is Mytheus's opening: the no-silent-commit contract turns every candidate into a confirm-or-reject event, precisely the evaluation signal the authors say the field lacks.
\end{takeaway}

\subsection{Limitations and what it leaves open}
\begin{itemize}[leftmargin=*,itemsep=1pt,topsep=2pt]
\item Every experiment runs on a static offline corpus (HumanEval, HotPotQA, NovelQA, WebArena) with a machine-checkable answer: no agent is evaluated on a live stream, none must act before its input is complete, and no task has a correct output that differs per user.
\item Cost is the only second axis; latency is named once as something the formulation should generalize to but never measured or optimized, so the framework gives no guidance under a hard deadline, and all three winning baselines work by adding serial model calls.
\item Only single-agent architectures are evaluated, with no agent-to-agent messaging, hence no accounting of the tokens, latency or failure modes that inter-agent traffic contributes to the cost curve.
\item The holdout ladder is prescribed but never validated: they count how many of 17 benchmarks have appropriate holdouts, yet never construct a compliant held-out set and show that rankings change under it, so the overfitting they warn about stays unquantified outside the qualitative STeP argument.
\item The cost model is single-tenant list-price arithmetic over independent runs; semantic and prompt caching, cross-user batching, per-tenant quotas and shared GPU serving of retrieval models sit outside it, and the optimizer tunes one prompt for one task rather than allocating a bounded budget across tenants.
\end{itemize}
